\documentclass[review,authoryear,12pt]{elsarticle}

\usepackage[T1]{fontenc}
\usepackage[utf8]{inputenc}

\begin{document}

\begin{frontmatter}

\title{Credible on Paper? Green Banking Policy, Access to Finance, and Institutional Capacity in the Greening of Firms Worldwide}

\author[buv]{Quang-Vinh Dang\corref{cor1}}
\ead{vinh.dq4@buv.edu.vn}

\author[bav]{Thi-Hong-Hanh Nguyen}
\ead{hanhnth1@hvnh.edu.vn}

\cortext[cor1]{Corresponding author.}

\address[buv]{British University Vietnam, Hung Yen, Vietnam}
\address[bav]{Banking Academy Vietnam, Hanoi, Vietnam}

\begin{abstract}
Green banking policies, coordinated through the Sustainable Banking and Finance Network (SBFN) across sixty-plus emerging economies, assume that supervisory guidelines directing banks to price environmental risk strengthen the link between firms' credit access and green investment. Using two independent World Bank Enterprise Survey samples (28{,}042 firms, 41 economies; 89{,}797 firms, 162 economies) and four estimators---pooled logit, Bayesian hierarchical, causal forest, and country fixed effects---we find bank credit access robustly predicts green-practice adoption (13 percentage points), but neither SBFN adoption nor regulatory quality moderates this relationship. The estimates are precise, not merely insignificant: the causal forest's credit effect is 0.072 for non-adopter firms versus 0.076 for adopter firms, and is nearly identical across regulatory-quality terciles, ruling out moderation of policy-relevant size. A country-level staggered event-study (217 economies) corroborates this at the macro level. Heterogeneity instead concentrates at firm size, relocating rather than resolving which institutional margin matters.
\end{abstract}

\begin{keyword}
green banking policy \sep access to finance \sep institutional complementarity \sep firm-level green investment \sep causal forest \sep hierarchical Bayesian model
\JEL G21 \sep O16 \sep Q56 \sep D22 \sep C11
\end{keyword}

\end{frontmatter}

\vspace{1cm}
\noindent\textbf{Corresponding author:} Quang-Vinh Dang, British University Vietnam, Hung Yen, Vietnam. Email: \texttt{vinh.dq4@buv.edu.vn}.

\vspace{0.5cm}
\noindent\textbf{ORCID:} Quang-Vinh Dang: 0000-0002-3877-8024. Thi-Hong-Hanh Nguyen: 0009-0006-7821-7404.

\vspace{0.5cm}
\noindent\textbf{Acknowledgements:} We thank the Enterprise Analysis Unit of the Development Economics Global Indicators Department of the World Bank Group for making the Enterprise Surveys data available.

\vspace{0.5cm}
\noindent\textbf{Declaration of interest:} Declarations of interest: none.

\vspace{0.5cm}
\noindent\textbf{Funding:} This research did not receive any specific grant from funding agencies in the public, commercial, or not-for-profit sectors.

\vspace{0.5cm}
\noindent\textbf{CRediT authorship contribution statement:} \textbf{Quang-Vinh Dang:} Conceptualization, Methodology, Software, Formal analysis, Data curation, Writing -- original draft. \textbf{Thi-Hong-Hanh Nguyen:} Conceptualization, Investigation, Validation, Writing -- review \& editing.

\end{document}
